The numerical ratio between existing marriages and widows, assumed for the professors' widows' fund in \textsc{Kritter}'s opinion to be \(2:1\), may differ greatly between societies drawn from different walks of life. If, to begin with, we compare a widows' fund society such as ours with all the existing marriages and widows in an entire country, considering only the general determining factor stated above [p. 140, top], we readily see that more widows must be counted for a given number of marriages in the former than in the latter. Among the great mass of a country's inhabitants, marriages take place at an earlier age on average, and the average difference between the ages of husband and wife is smaller than among university professors\footnote{For our widows' fund I have age data for only 62 married couples, from which the average difference is 9 years. For the whole population of a country one could hardly assume an average difference even half as large.}. Moreover, some members enter our widows' fund already married, whereas in the lists for an entire state all married couples are counted from the time of their marriage. But the effect of these two causes is small compared with the influence of another circumstance left aside in the abstract treatment above, pp. [140 ff.]: widows may leave the count through remarriage as well as through death. In a widows' fund where a widow who remarries loses all further claim to a pension, remarriage among widows tends to be rare. Our fund, in particular, records only a single such case in the whole period of its existence. For an entire country, however, it is assumed on the basis of experience that about a thirtieth of all widows leave the count each year through remarriage, and this substantially changes the ratio of existing marriages to widows. Finally, in a country whose population has been increasing for some time, this increase also has a substantial effect on the ratio of \textit{coexisting} marriages and widows: there will then be more existing marriages per widow than there would be without that circumstance.

It is usually stated as an empirical rule that for an entire country four existing marriages are to be counted per widow. The figures given by \textsc{Quetelet} for Belgium agree almost exactly with this ratio. For the whole Kingdom of Hanover, the counts of 1833—1842 give a somewhat smaller figure, namely 3,74 existing marriages per widow. But if the individual regions are distinguished, very large differences appear; for example,
\SourceLeaf{150}
\begin{center}
\begin{tabular}{l@{\;}r}
in the Landdrostei of Stade .\,.\,.\,.\,.\,.\,.\,.\,.\,.\,.\,. & 4,21\\
in the Landdrostei of Aurich .\,.\,.\,.\,.\,.\,.\,.\,.\,.\,. & 3,28\\
in the Berghauptmannschaft of Clausthal .\,.\,.\,. & 2,56
\end{tabular}
\end{center}
\noindent existing marriages per widow.

{\rightskip=0pt plus 1em
In the second half of the previous century, a number of widows' funds were seen to collapse because they had been founded on the assumption (entirely false in light of the remarks above)\footnote{Some, admittedly, also because of \textit{even} grosser errors.} that the ratio of existing marriages to widows found from counts for an entire country could also be regarded as authoritative for widows' funds. \textsc{Kritter} was one of those who never tired of opposing such false principles. But whenever he makes definite statements of his own in his writings, he \textit{either} merely says that \textit{at the very least} one widow must be counted per two marriages (e.g. \textsc{Kritter}'s Vorstellung des bisherigen Erfolgs u. s. w., p. 22); \textit{or} he expressly stipulates (Prufung eines Aufsatzes in der Berliner Monatsschrift, p. 28) that the ratio \(2:1\) may be assumed only if a sixth of all widows remarry and their pensions thereby cease; \textit{or} he gives figures yielding substantially \textit{more} widows than half the number of existing marriages; for example:\par}

\begin{quote}
in the work just cited, p. 27, a calculation gives, for 3293 married couples, 1831 widows,

Prüfung einer kleinen Schrift u. s. w., p. 17, where, for 100 husbands, 60 widows are counted.
\end{quote}
\noindent This same ratio is given in
\begin{quote}
Sammlung dreier Aufsätze über die Calenbergische, Preussische und Dänische Witwenversorgungsanstalten, pp. 30—35.

Sammlung wichtiger Erfahrungen u. s. w., p. 35, where \textsc{Kritter} counts, for 450 husbands, 280 widows.
\end{quote}
\noindent\hspace*{\parindent}Moreover, \textsc{Kritter} nowhere based his figures on actual direct experience (adequate experience of this kind could hardly have been obtained from Germany at that time), but on mortality tables and certain assumptions about the average age relations of those entering the society. These assumptions are in any case no less favorable than the circumstances actually found in our professors' widows' fund.

{\rightskip=0pt plus 1em
I find actual experience from a large group, probably roughly comparable to our widows' fund, in \textsc{Price}, Observations on reversionary payments, pp. 79. 269. 276 (the third edition of this work, of 1773). There, for all ministers and professors in Scotland, an average of 17 years gives 667 existing marriages and 380 widows, thus very nearly in the ratio of 7 to 4. It is also noted that these occupational classes there began to receive the income from their posts at an average age of 27 years, thus certainly earlier than may be assumed on average for professors at German universities. Remarriage among those widows was also less rare than among the widows of Göttingen professors.\par}


Even apart from its small size, the normal ratio applicable to our society cannot be derived from its own direct experience, because the changes in the number of existing marriages within it are unknown. In the whole hundred years of the institution's existence, this number is known for only two points in time: the present moment and, according to P.'s statement (above, p. 132), the time of the deliberations
\SourceLeaf{151}
of 1794. Had it been regularly recorded, for all 204 professors who have so far joined the widows' fund, whether they were married and during which part of their membership, together with the exact ages of both themselves and their wives, this could be used with the mortality tables to determine the ratio in question very accurately by an indirect method. But to my knowledge none of this has been done. The more regrettable this is now, the more confidently one may hope that appropriate measures, taken at least from now on, will spare the person who gives an opinion after another 100 years the same cause for complaint.

